\begin{enumerate}

\item La figura establece qué partes pertenecen a la sección, la dirección de y y z y las dimensiones en milímetros.

\item El área total suma las áreas macizas y resta los huecos; los primeros momentos proporcionan las coordenadas del centroide.

\item Calcular I\_z por suma de las inercias propias de las piezas más A\_i por el cuadrado de cada distancia a G (Steiner).

\item El verificador geométrico de Studium contrasta las inercias por otra vía: integración directa sobre los contornos.

\item La tensión normal combina N/A y la flexión alrededor de los ejes centrales; si I\_yz no es cero se conserva en la relación general.

\item Determinar las fibras de tensión extrema evaluando la distribución lineal en los vértices y representar las zonas de compresión y tracción.

\end{enumerate}

\noindent\textbf{Operaciones y resultados intermedios.}

\begin{enumerate}

\item Área algebraica de la sección, descontando los huecos.

\[\fitmath{y_{1}=\left(60-0\right)\cdot \left(120-0\right)\approx 7200\,\mathrm{mm^{2}}}\]

\item Centroide en el eje vertical mediante primeros momentos.

\[\fitmath{y_{2}=\frac{\frac{\left(60-0\right)\cdot \left(120-0\right)\cdot \left(0+120\right)}{2}}{\left(60-0\right)\cdot \left(120-0\right)}\approx 60\,\mathrm{mm}}\]

\item Momento de inercia I\_z mediante el teorema de Steiner.

\[\fitmath{y_{3}=\frac{\left(60-0\right)\cdot {\left(120-0\right)}^{3}}{12}+\left(60-0\right)\cdot \left(120-0\right)\cdot {\left(\frac{0+120}{2}-60\right)}^{2}\approx 8640000\,\mathrm{mm^{4}}}\]

\item Tensión extrema: axil y flexión con el convenio declarado.

\[\fitmath{y_{4}=\frac{250\cdot \sqrt{3}}{9}+\frac{500}{9}\approx 103.6681\,\mathrm{MPa}}\]

\end{enumerate}
